\documentclass[aspectratio=169,11pt]{beamer}

\usepackage[utf8]{inputenc}
\usepackage[T1]{fontenc}
\usepackage[french]{babel}
\usepackage{lmodern}
\usepackage{xcolor}
\usepackage{tikz}

% ---- Thème épuré, accent orange de l'application ----
\usetheme{default}
\setbeamertemplate{navigation symbols}{}
\setbeamertemplate{caption}[numbered]

\definecolor{accent}{HTML}{D97757}
\definecolor{dark}{HTML}{243857}
\definecolor{soft}{HTML}{6B7280}

\setbeamercolor{structure}{fg=accent}
\setbeamercolor{frametitle}{fg=dark}
\setbeamercolor{title}{fg=dark}
\setbeamercolor{itemize item}{fg=accent}
\setbeamercolor{itemize subitem}{fg=accent}
\setbeamerfont{frametitle}{size=\Large,series=\bfseries}
\setbeamertemplate{itemize item}{\textbullet}

% Barre de pied de page minimaliste
\setbeamertemplate{footline}{%
  \hbox{\begin{beamercolorbox}[wd=\paperwidth,ht=2.6ex,dp=1.2ex,leftskip=1em,rightskip=1em]{}%
  \color{soft}\small Voyage Assistant --- SAE2 BUT3 \hfill \insertframenumber/\inserttotalframenumber
  \end{beamercolorbox}}%
}

\newcommand{\code}[1]{\texttt{\small #1}}
\newcommand{\hl}[1]{\textbf{\color{accent}#1}}

\title{\textbf{Voyage Assistant}}
\subtitle{Assistant IA conversationnel pour le transport}
\date{SAE2 BUT3 --- Promotion 2025--2026}

\begin{document}

% ---- Titre ----
\begin{frame}[plain]
  \centering
  \vspace{0.6cm}
  {\color{accent}\rule{0.85\linewidth}{1.5pt}}\\[0.5cm]
  {\Huge\bfseries\color{dark} Voyage Assistant}\\[0.3cm]
  {\large\color{soft} Assistant IA conversationnel pour le transport}\\[0.4cm]
  {\color{accent}\rule{0.85\linewidth}{1.5pt}}\\[0.7cm]
  {\footnotesize
   Ako Christian \quad\textbullet\quad Carl Similien \quad\textbullet\quad
   Noé Cervera \quad\textbullet\quad Dhanoush Kessavane}\\[0.5cm]
  {\small Encadrement : \textbf{Bilal Faye} (LIPN, CNRS UMR 7030)}\\[0.5cm]
  {\scriptsize\color{soft} optimalako-voyage-assistant.hf.space}
\end{frame}

% ---- Équipe ----
\begin{frame}{L'équipe et les rôles}
  \begin{itemize}
    \item \hl{Ako Christian} --- IA conversationnelle \& RAG, déploiement
    \item \hl{Carl Similien} --- Back-end \& intégration (API, base de données)
    \item \hl{Noé Cervera} --- Modèles \& évaluation (embeddings, reranker, NER)
    \item \hl{Dhanoush Kessavane} --- Front-end, design \& documentation
  \end{itemize}
  \vfill
  \begin{center}\small\color{soft}
    Encadrement : Bilal Faye (LIPN, CNRS UMR 7030)
  \end{center}
\end{frame}

% ---- Contexte ----
\begin{frame}{Contexte et objectif}
  \textbf{Domaine choisi : le transport de voyageurs} (avion, train, bateau, bus).
  \vspace{0.4cm}
  \begin{itemize}
    \item Une application \hl{full-stack} de réservation de trajets.
    \item Un \hl{assistant IA} qui répond aux questions \emph{et} agit réellement
          (modification, réclamation, annulation).
    \item Idée directrice : \hl{un LLM seul ne suffit pas}. Il invente et ignore les
          données récentes ou personnelles.
    \item Notre réponse : coupler le LLM à des \hl{sources de vérité} via le RAG.
  \end{itemize}
\end{frame}

% ---- Démo ----
\begin{frame}{L'application est en ligne}
  \begin{center}
    \vspace{0.3cm}
    {\Large Déployée et accessible publiquement :}\\[0.4cm]
    {\large\color{accent}\bfseries https://optimalako-voyage-assistant.hf.space}\\[0.8cm]
    \begin{minipage}{0.8\linewidth}\centering\color{soft}
      Recherche parmi plus de \textbf{182\,000} liaisons réelles, réservation avec
      billet PDF et e-mail, compte Google ou e-mail, et l'assistant conversationnel.
    \end{minipage}
  \end{center}
\end{frame}

% ---- Architecture ----
\begin{frame}{Architecture technique}
  \begin{columns}[T]
    \begin{column}{0.5\linewidth}
      \textbf{\color{dark}Frontend}
      \begin{itemize}\small
        \item React 18 (sans build)
        \item Thèmes clair / sombre
        \item Voix : Web Speech API
      \end{itemize}
      \vspace{0.3cm}
      \textbf{\color{dark}Backend}
      \begin{itemize}\small
        \item FastAPI + Pydantic
        \item SQLAlchemy (ORM)
        \item Services IA (LLM, RAG, NER)
      \end{itemize}
    \end{column}
    \begin{column}{0.5\linewidth}
      \textbf{\color{dark}Base de données}
      \begin{itemize}\small
        \item PostgreSQL (Neon) en prod
        \item SQLite en local
      \end{itemize}
      \vspace{0.3cm}
      \textbf{\color{dark}Infrastructure}
      \begin{itemize}\small
        \item Hugging Face Spaces (Docker)
        \item Un seul service : API + RAG + front
      \end{itemize}
    \end{column}
  \end{columns}
\end{frame}

% ---- RAG ----
\begin{frame}{Qu'est-ce que le RAG et la vectorisation ?}
  \hl{RAG} = \emph{Retrieval-Augmented Generation} : on \hl{récupère l'information
  pertinente avant de générer} la réponse, puis on l'injecte dans le prompt.
  \vspace{0.4cm}
  \begin{itemize}
    \item \hl{Vectorisation} : chaque texte devient un vecteur de nombres
          (384 dimensions) qui encode son \emph{sens}.
    \item Deux textes de sens proche ont des vecteurs proches \hl{même sans mots
          communs}.
    \item Exemple : « je veux annuler » est relié à l'article CGV
          « conditions de résiliation ».
  \end{itemize}
  \vfill
  \begin{center}\small\color{soft}
    Analogie : un examen « à livre ouvert » plutôt que de réciter de mémoire.
  \end{center}
\end{frame}

% ---- Pipeline RAG ----
\begin{frame}{La chaîne RAG en deux étapes}
  \begin{enumerate}
    \item \hl{Indexation} : documents découpés en \emph{chunks}, vectorisés par
          \code{gte-small} (en local), stockés dans \hl{FAISS}.
    \item \hl{Recherche} en deux temps :
      \begin{itemize}\small
        \item \textbf{Bi-encoder} : rapide, remonte 10 candidats.
        \item \textbf{Cross-encoder} (\emph{reranking}) : précis, garde les 3 meilleurs.
      \end{itemize}
  \end{enumerate}
  \vspace{0.3cm}
  Plus un \hl{NER} (\code{camembert-ner}) qui extrait prénom et nom d'une phrase libre.
  \vfill
  \begin{center}\small\color{soft}
    On obtient la précision du cross-encoder sans son coût sur tout le corpus.
  \end{center}
\end{frame}

% ---- 3 sources ----
\begin{frame}{Les trois sources de connaissance}
  Selon la question, l'assistant interroge la bonne source :
  \vspace{0.3cm}
  \begin{itemize}
    \item \hl{RAG documentaire} (FAISS) --- les règles : annulation, bagages\dots
    \item \hl{Base de données} --- les questions personnelles.\\
          \small\emph{« Combien de temps va durer mon vol ? »} $\rightarrow$ lecture du
          billet, calcul de la durée, réponse personnalisée.
    \item \hl{APIs temps réel + web} --- les données du jour : Open-Meteo (météo,
          heure), Frankfurter (change), DuckDuckGo.
  \end{itemize}
  \vfill
  \begin{center}\small\color{soft}
    Les infos stables (capitale, monnaie) ne sont pas appelées : le LLM les connaît.
  \end{center}
\end{frame}

% ---- Actions réelles ----
\begin{frame}{Des actions réelles, sans hallucination}
  Modifier, annuler, réclamer ne sont \hl{jamais} confiés au LLM.\\[0.3cm]
  Une \hl{machine à états} déterministe pilote le parcours :
  \vspace{0.3cm}
  \begin{center}\small
    intention $\rightarrow$ vérification d'identité $\rightarrow$
    \textbf{exécution réelle en base} $\rightarrow$ e-mail + billet PDF
  \end{center}
  \vspace{0.3cm}
  \begin{itemize}
    \item Le LLM \hl{rédige} seulement ; il \hl{n'exécute} jamais.
    \item Cette séparation élimine les fausses confirmations.
  \end{itemize}
\end{frame}

% ---- Sécurité ----
\begin{frame}{Sécurité}
  \begin{itemize}
    \item Garde-fous : insultes, \hl{injections de prompt}, sujets hors-périmètre.
    \item Anti-extraction : pas d'info sans identité vérifiée.
    \item SQL paramétré (anti-injection), échappement HTML de React (anti-XSS).
    \item Triple vérification d'identité (numéro, nom, prénom, date de naissance).
    \item \hl{Aucun secret dans le code} : variables d'environnement uniquement.
  \end{itemize}
\end{frame}

% ---- Déploiement ----
\begin{frame}{Déploiement : les choix}
  \begin{itemize}
    \item Le RAG charge \code{torch} + 2 modèles $\rightarrow$ \hl{700 Mo à 1 Go de RAM}.
    \item Render / Railway gratuits (512 Mo) $\rightarrow$ \hl{dépassement mémoire}.
    \item Choix : \hl{Hugging Face Spaces} (16 Go gratuits, conçu pour le ML).
    \item Un seul service Docker : API + RAG + front (même origine, pas de CORS).
    \item Base \hl{Neon} (PostgreSQL) co-localisée pour minimiser la latence.
  \end{itemize}
  \vfill
  \begin{center}\small\color{soft}
    Image Docker en torch CPU (200 Mo) ; migration des 182\,k routes vers Neon.
  \end{center}
\end{frame}

% ---- Difficultés ----
\begin{frame}{Difficultés rencontrées et solutions}
  \begin{itemize}\small
    \item DuckDuckGo limité $\rightarrow$ librairie \code{ddgs} (rotation multi-backend).
    \item Pas de date de naissance via Google $\rightarrow$ formulaire = source de vérité.
    \item Hallucinations d'action $\rightarrow$ machine à états séparée du LLM.
    \item Recherche d'escales à 15 s $\rightarrow$ index SQL + cache (0,34 s).
    \item RAG « classique » $\rightarrow$ LangChain + FAISS + reranking + NER.
    \item \code{torch} CUDA de 1,2 Go $\rightarrow$ version CPU dans le Docker.
    \item 512 Mo insuffisants $\rightarrow$ Hugging Face Spaces (16 Go).
  \end{itemize}
\end{frame}

% ---- Cahier des charges ----
\begin{frame}{Réponse au cahier des charges}
  \begin{itemize}\small
    \item Domaine spécialisé : \hl{transport}.
    \item Full-stack : \hl{FastAPI + React + PostgreSQL}.
    \item LLM : \hl{Gemini 2.5 Flash}.
    \item RAG / vectorisation : \hl{LangChain + FAISS + gte-small + reranker}.
    \item Base de session : \hl{conversations persistées}.
    \item Voix (option) : \hl{Web Speech API}.
    \item Recherche web (bonus) : \hl{DuckDuckGo + APIs temps réel}.
    \item Déploiement : \hl{en ligne sur Hugging Face Spaces + Neon}.
  \end{itemize}
\end{frame}

% ---- Conclusion ----
\begin{frame}{Conclusion et perspectives}
  \begin{itemize}
    \item Toutes les exigences de la SAE2 sont \hl{remplies}, déploiement inclus.
    \item L'IA s'appuie sur des \hl{sources de vérité} pour rester factuelle.
    \item Perspectives : plus de documents métier, notifications temps réel,
          intégration continue.
  \end{itemize}
  \vfill
  \begin{center}
    {\Large\color{dark}\bfseries Merci de votre attention}\\[0.3cm]
    {\small\color{soft} Questions ?}
  \end{center}
\end{frame}

\end{document}
